\documentclass[12pt]{article}
\usepackage{graphicx}
\usepackage{hyperref}
\usepackage{subcaption}
\usepackage{xepersian}
\usepackage{listings}
\usepackage{xcolor}
\usepackage{float}



\lstset{
  language=C,
  basicstyle=\ttfamily\footnotesize,
  keywordstyle=\color{blue},
  stringstyle=\color{orange},
  commentstyle=\color{gray},
  numbers=left,
  numberstyle=\tiny\color{gray},
  frame=single,
  breaklines=true,
  captionpos=b
}


\setcounter{secnumdepth}{0}
\settextfont[Kashida]{XB Niloofar}
\setkeys{Gin}{width =\textwidth }

\title {پروژه سوم سیستم های عامل}
\author {دلارام روحانی 40231026 \\
        امیرعباس انتظاری 40231006}



\begin{document}
\maketitle
\tableofcontents
\newpage

\section{\lr{Swapping}}

\subsection{مقدمه}
سیستم‌عامل \lr{xv6} به‌صورت پیش‌فرض تنها از حافظه‌ی اصلی (\lr{RAM}) برای نگهداری صفحات استفاده می‌کند و در صورت پر شدن حافظه، تخصیص حافظه با شکست مواجه می‌شود. هدف فاز اول این پروژه، پیاده‌سازی مکانیزم \lr{Swap Memory} است تا سیستم‌عامل بتواند در شرایط کمبود حافظه، صفحات کم‌استفاده را به دیسک منتقل کرده و در صورت نیاز مجدداً آن‌ها را به حافظه بازگرداند.

در این فاز، تمرکز اصلی بر طراحی و پیاده‌سازی سه بخش زیر است:
\begin{itemize}
  \item ایجاد فرآیند کرنلی برای مدیریت \lr{Swap} (daemon)
  \item پیاده‌سازی \lr{Swap-out} (خارج کردن صفحه از حافظه)
  \item پیاده‌سازی \lr{Swap-in} (بازگرداندن صفحه به حافظه هنگام \lr{Page Fault})
\end{itemize}

\subsection{نمای کلی طراحی}
ایده‌ی اصلی این فاز مشابه سیستم‌های عامل واقعی است:
\begin{itemize}
  \item حافظه‌ی اصلی نقش یک کش را دارد.
  \item دیسک نقش فضای پشتیبان \lr{(Backing Store)} را ایفا می‌کند.
  \item جدول صفحات مشخص می‌کند که هر صفحه در حال حاضر در \lr{RAM} است یا روی دیسک.
\end{itemize}

طبق صورت پروژه، برای جلوگیری از پیچیدگی‌های کرنل و مشکلات ناشی از عملیات فایل در کرنل، خواندن/نوشتن فایل‌های \lr{swap} در \textbf{فضای کاربر} انجام می‌شود و کرنل تنها مسئول هماهنگی، انتخاب صفحه قربانی، تغییر \lr{PTE} و آزادسازی/بازگردانی حافظه است.

\subsection{ایجاد فرآیند کرنلی \texttt{swapd}}
برای مدیریت عملیات \lr{Swap-out} یک فرآیند کرنلی دائمی (daemon) به نام \texttt{swapd} ایجاد شده است. این فرآیند مانند یک پردازه‌ی عادی توسط زمان‌بند اجرا می‌شود، اما هرگز وارد فضای کاربر نمی‌شود و تمام عمر خود را در کرنل سپری می‌کند.

\subsubsection{توسعه ساختار پردازه}
در فایل \texttt{proc.h} به ساختار \texttt{proc} دو فیلد جدید اضافه شد:
\begin{LTR}
\begin{lstlisting}
int is_kproc;
void (*kentry)(void);
\end{lstlisting}
\end{LTR}

فیلد \texttt{is\_kproc} مشخص می‌کند پردازه کرنلی است یا نه، و \texttt{kentry} اشاره‌گری به تابعی است که نقطه‌ی شروع اجرای پردازه در کرنل محسوب می‌شود.

\subsubsection{ساخت پردازه کرنلی}
در فایل \texttt{proc.c} تابع \texttt{create\_kernel\_process} پیاده‌سازی شد. این تابع یک \lr{proc} جدید می‌سازد، آن را به‌عنوان پردازه‌ی کرنلی علامت‌گذاری می‌کند، زمینه‌ی اجرایی آن را طوری تنظیم می‌کند که به‌جای \texttt{forkret} وارد تابع شروع کرنلی شود، سپس آن را \lr{RUNNABLE} کرده و وارد صف زمان‌بند می‌کند.

همچنین تابع \texttt{kproc\_start} به‌عنوان نقطه‌ی ورود مشترک پردازه‌های کرنلی اضافه شد: ابتدا قفل پردازه آزاد می‌شود، وقفه‌ها فعال می‌شوند، سپس \texttt{kentry} اجرا می‌گردد و اگر تابع بازگردد، پردازه با \texttt{yield} در حلقه‌ی بی‌نهایت باقی می‌ماند تا هیچ‌گاه «از کرنل خارج نشود».

\subsubsection{رفتار \texttt{swapd}}
تابع \texttt{swapd} به صورت حلقه‌ی بی‌نهایت اجرا می‌شود:
\begin{itemize}
  \item اگر درخواستی برای آزادسازی حافظه وجود نداشته باشد، \texttt{sleep} می‌کند.
  \item با ورود درخواست، یک صفحه قربانی انتخاب و \lr{Swap-out} می‌شود.
  \item پس از آزاد شدن حداقل یک صفحه فیزیکی، تمام پردازه‌های منتظر حافظه بیدار می‌شوند.
\end{itemize}

\subsection{شروع \lr{Swap-out}: نقطه‌ی ورود در \texttt{uvmalloc}}
شروع \lr{Swap-out} زمانی رخ می‌دهد که تخصیص حافظه‌ی فیزیکی شکست بخورد.

در فایل \texttt{vm.c}، در تابع \texttt{uvmalloc}، به جای آنکه با شکست \texttt{kalloc()} عملیات تخصیص متوقف شود، پردازه وارد حالت انتظار می‌شود تا \texttt{swapd} یک صفحه آزاد کند:
\begin{LTR}
\begin{lstlisting}
while((mem = kalloc()) == 0){
  swap_wait_for_free_page();
}
\end{lstlisting}
\end{LTR}

به این ترتیب، تخصیص حافظه در سطح \lr{user} شکست نمی‌خورد، بلکه کرنل به صورت کنترل‌شده فضا ایجاد می‌کند.

\subsection{صف درخواست‌های \lr{Swap} و مفهوم \texttt{chan}}
در فایل \texttt{swap.c} یک صف حلقوی (\lr{ring buffer}) برای درخواست‌ها پیاده‌سازی شد. برای هماهنگی بین پردازه‌ها و \texttt{swapd} از مکانیزم \texttt{sleep/wakeup} در xv6 استفاده شده است.

در xv6، \texttt{sleep(chan, lock)} پردازه را روی «کانال» مشخص‌شده می‌خواباند و \texttt{wakeup(chan)} تمام پردازه‌های خوابیده روی آن کانال را بیدار می‌کند. \texttt{chan} صرفاً یک آدرس (\texttt{void*}) است که به عنوان کلید همگام‌سازی استفاده می‌شود.

دو کانال اصلی در این فاز:
\begin{itemize}
  \item \texttt{swapq\_chan}: کانالی که \texttt{swapd} هنگام خالی بودن صف روی آن می‌خوابد.
  \item \texttt{swap\_wait\_chan}: کانالی که پردازه‌های منتظر آزاد شدن حافظه روی آن می‌خوابند.
\end{itemize}

تابع \texttt{swap\_wait\_for\_free\_page} درخواست را در صف قرار می‌دهد، \texttt{swapd} را بیدار می‌کند، سپس خود پردازه را تا آزاد شدن حداقل یک صفحه فیزیکی روی \texttt{swap\_wait\_chan} می‌خواباند.

\subsection{انتخاب صفحه‌ی قربانی: تقریب \lr{LRU} با \texttt{PTE\_A}}
برای انتخاب صفحه مناسب جهت \lr{Swap-out} از یک تقریب \lr{LRU} به سبک \lr{Second-Chance} استفاده شد.

در معماری \lr{RISC-V} سخت‌افزار هنگام دسترسی به یک صفحه، بیت \texttt{Accessed} را در \lr{PTE} فعال می‌کند (بیت \texttt{PTE\_A}). الگوریتم انتخاب قربانی به صورت دو گذر اجرا می‌شود:
\begin{itemize}
  \item \textbf{گذر اول:} اگر صفحه \texttt{PTE\_A=1} داشته باشد، بیت آن پاک می‌شود تا «فرصت دوم» داده شود.
  \item \textbf{گذر دوم:} اولین صفحه‌ی قابل تخلیه که \texttt{PTE\_A=0} باشد به عنوان قربانی انتخاب می‌شود.
\end{itemize}

همچنین صفحات غیرقابل تخلیه نادیده گرفته می‌شوند (مثلاً صفحات کرنل، صفحات غیرمعتبر، یا صفحاتی که قبلاً \lr{swap} شده‌اند).

\subsection{بیت \texttt{PTE\_SWP} و دلیل نیاز به آن}
صرف صفر کردن \texttt{PTE\_V} کافی نیست؛ زیرا یک \lr{PTE} نامعتبر می‌تواند دو معنی کاملاً متفاوت داشته باشد:
\begin{itemize}
  \item صفحه واقعاً نگاشت نشده است (دسترسی غیرمجاز)
  \item صفحه نگاشت شده بوده اما به دیسک منتقل شده است (قابل بازیابی)
\end{itemize}

برای تفکیک این دو حالت، یک بیت نرم‌افزاری جدید به \texttt{riscv.h} اضافه شد:
\begin{LTR}
\begin{lstlisting}
#define PTE_SWP (1L << ...)
\end{lstlisting}
\end{LTR}

معانی حالت‌ها:
\begin{itemize}
  \item \texttt{PTE\_V=1}: صفحه در \lr{RAM} قرار دارد.
  \item \texttt{PTE\_V=0} و \texttt{PTE\_SWP=1}: صفحه \lr{swap} شده و قابل بازیابی است.
  \item \texttt{PTE\_V=0} و \texttt{PTE\_SWP=0}: دسترسی نامعتبر (برنامه باید کشته شود).
\end{itemize}

\subsection{نوشتن صفحه روی دیسک در فضای کاربر (\texttt{swapper})}
طبق صورت پروژه، فایل \lr{swap} باید به‌ازای هر صفحه ساخته شود و فرمت نام‌گذاری آن:
\[
\texttt{pid\_[L2,L1,L0].swp}
\]
است تا بازیابی صفحه در آینده امکان‌پذیر باشد.

به دلیل ممنوعیت/عدم توصیه عملیات فایل در کرنل، یک برنامه‌ی کاربری به نام \texttt{swapper} پیاده‌سازی شد. کرنل و این برنامه از طریق یک \lr{mailbox} مشترک و چند \lr{syscall} با یکدیگر ارتباط برقرار می‌کنند:
\begin{itemize}
  \item کرنل هنگام \lr{Swap-out} داده صفحه را در بافر مشترک قرار داده و یک \texttt{task} ثبت می‌کند.
  \item \texttt{swapper} با \texttt{swap\_fetch} درخواست را دریافت می‌کند.
  \item بسته به نوع عملیات (\lr{WRITE/READ}) فایل \texttt{.swp} را می‌نویسد یا می‌خواند.
  \item سپس با \texttt{swap\_complete} (و در حالت \lr{READ} با نسخه‌ای که داده را برمی‌گرداند) نتیجه را به کرنل اعلام می‌کند.
\end{itemize}

\subsection{مراحل \lr{Swap-out} در \texttt{swapout\_one\_page}}
عملیات کامل \lr{Swap-out} به صورت زیر است:
\begin{enumerate}
  \item انتخاب صفحه قربانی با الگوریتم \lr{LRU} تقریبی.
  \item استخراج آدرس فیزیکی صفحه از \lr{PTE}.
  \item ارسال درخواست \lr{WRITE} به برنامه‌ی \texttt{swapper} و انتظار برای اتمام نوشتن.
  \item به‌روزرسانی \lr{PTE}: صفر کردن \texttt{PTE\_V} و فعال‌کردن \texttt{PTE\_SWP}.
  \item آزاد کردن صفحه فیزیکی با \texttt{kfree} تا حافظه در دسترس تخصیص‌های بعدی قرار گیرد.
\end{enumerate}

پس از آزاد شدن حداقل یک صفحه، \texttt{swapd} پردازه‌های منتظر روی \texttt{swap\_wait\_chan} را بیدار می‌کند تا تخصیص حافظه ادامه یابد.

\subsection{\lr{Swap-in}: بازگردانی صفحه هنگام \lr{Page Fault}}
\lr{Swap-in} زمانی رخ می‌دهد که پردازه به صفحه‌ای دسترسی پیدا کند که در \lr{RAM} نیست اما \lr{swap} شده است.

\subsubsection{تشخیص در \texttt{usertrap}}
در فایل \texttt{trap.c} در تابع \texttt{usertrap}، هنگام وقوع \lr{page fault} (برای \lr{load/store} و در صورت نیاز \lr{instruction}) آدرس خطا از \texttt{stval} خوانده می‌شود و \lr{PTE} مربوط به آن \lr{VA} بررسی می‌گردد. اگر:
\[
(\texttt{PTE\_V} = 0) \land (\texttt{PTE\_SWP} = 1)
\]
باشد، یعنی صفحه \lr{swap} شده و باید \lr{swap-in} انجام شود.

\subsubsection{مراحل \texttt{swapin\_page}}
تابع \texttt{swapin\_page} مراحل زیر را انجام می‌دهد:
\begin{enumerate}
  \item تخصیص یک صفحه فیزیکی جدید (در صورت کمبود حافظه، دوباره از \texttt{swap\_wait\_for\_free\_page} استفاده می‌شود).
  \item ارسال درخواست \lr{READ} به \texttt{swapper} برای خواندن فایل \texttt{.swp}.
  \item کپی داده‌ی برگشتی به صفحه فیزیکی جدید.
  \item بازسازی \lr{PTE}: قرار دادن \lr{PA} جدید، فعال کردن \texttt{PTE\_V} و پاک کردن \texttt{PTE\_SWP}، و حفظ مجوزهای صفحه (\texttt{R/W/X/U}).
  \item اجرای \texttt{sfence\_vma} برای همگام‌سازی \lr{TLB}.
\end{enumerate}

پس از این مراحل، دستور خطادار مجدداً اجرا شده و برنامه بدون خطا ادامه می‌یابد.

\subsection{تست فاز اول (\texttt{swaptest})}
برای تست صحیح بودن \lr{Swap-out/Swap-in} یک برنامه‌ی کاربری \texttt{swaptest} نوشته شد که:
\begin{itemize}
  \item ۲۰ پردازه‌ی فرزند ایجاد می‌کند.
  \item هر فرزند ۱۰ بار حافظه‌ی ۴KB تخصیص می‌دهد.
  \item هر صفحه با الگوی مشخص پر می‌شود.
  \item در انتها همه‌ی صفحات دوباره بررسی می‌شوند تا فساد داده رخ نداده باشد.
\end{itemize}

برای فعال شدن سریع‌تر \lr{swap} طبق پیشنهاد صورت پروژه، مقدار \texttt{PHYSTOP} در فایل \texttt{memlayout.h} کاهش داده شد تا حافظه فیزیکی محدودتر شود و سیستم سریع‌تر به کمبود \lr{RAM} برسد.

نمایش پیام \texttt{swaptest: ALL OK} نشان‌دهنده‌ی موفقیت تست و در نتیجه تکمیل صحیح فاز اول است.

\subsection{جمع‌بندی}
در این فاز، یک سیستم \lr{Swap} کامل شامل \lr{Swap-out} و \lr{Swap-in} به xv6 اضافه شد. این پیاده‌سازی:
\begin{itemize}
  \item تحت فشار حافظه از شکست تخصیص جلوگیری می‌کند.
  \item با استفاده از daemon کرنلی \texttt{swapd} مدیریت \lr{eviction} را از مسیرهای حساس جدا می‌کند.
  \item با افزودن \texttt{PTE\_SWP} تفاوت بین صفحات نامعتبر و صفحات \lr{swap} شده را مشخص می‌کند.
  \item عملیات فایل را مطابق صورت پروژه در فضای کاربر انجام می‌دهد.
  \item با تست \texttt{swaptest} صحت عملکرد آن تأیید شده است.
\end{itemize}


\section{\lr{Namespace}ها}

\lr{Namespace}ها مکانیزمی در سیستم‌عامل هستند که برای ایزوله‌سازی منابع بین پردازه‌ها استفاده می‌شوند. با استفاده از \lr{Namespace}ها، هر پردازه می‌تواند دید مستقلی از بخش‌هایی از سیستم داشته باشد، به‌طوری که تغییرات آن روی سایر پردازه‌ها تأثیری نگذارد. این مفهوم پایه‌ی اصلی پیاده‌سازی \lr{Container}ها در سیستم‌عامل‌های مدرن است.

\subsection{\lr{PID Namespace}}
\lr{PID Namespace} باعث می‌شود هر پردازه شناسه‌ی پردازه‌ی (\lr{PID}) مخصوص به فضای نام خود را داشته باشد. در این حالت، یک پردازه ممکن است در فضای نام خودش \lr{PID} متفاوتی نسبت به سیستم اصلی داشته باشد. این قابلیت برای ایزوله‌سازی درخت پردازه‌ها و اجرای مستقل برنامه‌ها بسیار کاربردی است.

کد زیر مربوط به پیاده‌سازی \lr{PID Namespace} است. در این بخش، با استفاده از تابع
\texttt{\lr{pid\_namespace\_alloc}} یک فضای نام جدید برای شناسه‌های پردازه ایجاد می‌شود.
همچنین تابع \texttt{\lr{pid\_namespace\_get\_pid}} با استفاده از قفل، یک \lr{PID} یکتا
در داخل فضای نام تولید می‌کند تا از تداخل بین پردازه‌ها جلوگیری شود.


\begin{LTR}
\begin{lstlisting}
// تخصیص یک PID Namespace جدید
struct pid_namespace*
pid_namespace_alloc()
{
  struct pid_namespace *ns;
  
  ns = (struct pid_namespace *)kalloc();
  if(ns == 0)
    return 0;
  
  ns->refcount = 1;
  ns->next_pid = 1;
  initlock(&ns->lock, "pid_ns");
  
  return ns;
}

// دریافت بعدی PID از namespace
int
pid_namespace_get_pid(struct pid_namespace *ns)
{
  int pid;
  acquire(&ns->lock);
  pid = ns->next_pid;
  ns->next_pid = ns->next_pid + 1;
  release(&ns->lock);
  return pid;
}
\end{lstlisting}
\end{LTR}

\subsection{\lr{UTS Namespace}}
\lr{UTS Namespace} امکان جداسازی اطلاعات مربوط به سیستم مانند نام میزبان (\lr{Hostname}) را فراهم می‌کند. هر فضای نام \lr{UTS} می‌تواند \lr{hostname} مخصوص به خود را داشته باشد، بدون اینکه روی سایر پردازه‌ها یا سیستم اصلی تأثیر بگذارد.

کد زیر پیاده‌سازی \lr{UTS Namespace} را نشان می‌دهد. در این قسمت یک فضای نام
جدید برای اطلاعات سیستمی مانند \lr{Hostname} ایجاد می‌شود. با استفاده از شمارنده
ارجاع (\lr{refcnt}) و قفل، امکان اشتراک ایمن این فضا بین چند پردازه فراهم شده است.


\begin{LTR}
\begin{lstlisting}
// تخصیص یک UTS Namespace جدید
struct uts_namespace*
uts_namespace_alloc(void)
{
  struct uts_namespace *ns;
  
  ns = (struct uts_namespace *)kalloc();
  if(ns == 0)
    return 0;
  
  ns->refcnt = 1;
  memset(ns->hostname, 0, HOSTNAME_LEN);
  initlock(&ns->lock, "uts_ns");
  
  return ns;
}

void
uts_namespace_get(struct uts_namespace *ns)
{
  if(ns == 0)
    return;
  acquire(&ns->lock);
  ns->refcnt++;
  release(&ns->lock);
}
\end{lstlisting}
\end{LTR}

\subsection{\lr{IPC Namespace}}
\lr{IPC Namespace} برای ایزوله‌سازی منابع ارتباط بین‌پردازه‌ای مانند \lr{Semaphore}ها، \lr{Message Queue}ها و \lr{Shared Memory} استفاده می‌شود. با این کار، پردازه‌ها فقط می‌توانند به منابع \lr{IPC} داخل فضای نام خود دسترسی داشته باشند.

در این بخش، کد مربوط به \lr{IPC Namespace} ارائه شده است. این پیاده‌سازی
امکان ایزوله‌سازی منابع ارتباط بین‌پردازه‌ای را فراهم می‌کند. با تخصیص یک
\lr{IPC Namespace} جدید، پردازه‌ها فقط به منابع ارتباطی داخل فضای نام خود
دسترسی خواهند داشت.


\begin{LTR}
\begin{lstlisting}
// تخصیص یک IPC Namespace جدید
struct ipc_namespace*
ipc_namespace_alloc(void)
{
  struct ipc_namespace *ns;
  
  ns = (struct ipc_namespace *)kalloc();
  if(ns == 0)
    return 0;
  
  ns->refcnt = 1;
  initlock(&ns->lock, "ipc_ns");
  
  return ns;
}

void
ipc_namespace_get(struct ipc_namespace *ns)
{
  if(ns == 0)
    return;
  acquire(&ns->lock);
  ns->refcnt++;
  release(&ns->lock);
}
\end{lstlisting}
\end{LTR}

\subsection{\lr{Mount Namespace}}
\lr{Mount Namespace} باعث جداسازی فضای فایل سیستم می‌شود. هر پردازه می‌تواند \lr{mount point}های مخصوص به خود را داشته باشد و تغییرات مربوط به \lr{mount} یا \lr{unmount} فایل سیستم در یک \lr{namespace}، روی سایر \lr{namespace}ها اعمال نمی‌شود.

کد زیر پیاده‌سازی \lr{Mount Namespace} را نشان می‌دهد. در این فضا، هر پردازه
می‌تواند نمای مستقل خود از سیستم فایل را داشته باشد. ریشه‌ی فایل سیستم
در ساختار \lr{namespace} ذخیره شده و با استفاده از قفل، هماهنگی بین پردازه‌ها
حفظ می‌شود.

\begin{LTR}
\begin{lstlisting}
// تخصیص یک Mount Namespace جدید
struct mount_namespace*
mount_namespace_alloc(struct inode *root)
{
  struct mount_namespace *ns;
  
  ns = (struct mount_namespace *)kalloc();
  if(ns == 0)
    return 0;
  
  ns->refcnt = 1;
  ns->root = root;
  initlock(&ns->lock, "mount_ns");
  
  return ns;
}

void
mount_namespace_get(struct mount_namespace *ns)
{
  if(ns == 0)
    return;
  acquire(&ns->lock);
  ns->refcnt++;
  release(&ns->lock);
}
\end{lstlisting}
\end{LTR}

\section{ns\_test.c}

فایل \texttt{ns\_test.c} یک برنامه‌ی جامع برای تست و نمایش کارایی namespace های پیاده‌سازی شده است. این برنامه مختلف حالت‌های استفاده از namespace ها را بررسی می‌کند و نحوه ایزوله‌سازی منابع را نشان می‌دهد.

\textbf{توضیح برنامه:} این تست شامل پنج بخش اصلی است:

\begin{enumerate}
\item \textbf{\lr{Test 1}}: بررسی \texttt{\lr{PID Namespace}} و اینکه چگونه هر process یک \lr{PID} منحصربه‌فرد دریافت می‌کند.

\item \textbf{\lr{Test 2}}: تست \texttt{\lr{UTS Namespace}} برای تغییر و بررسی نام میزبان (\texttt{\lr{hostname}}).

\item \textbf{\lr{Test 3}}: استفاده از \texttt{\lr{unshare}} برای جداسازی \texttt{\lr{UTS Namespace}} و اینکه child process می‌تواند hostname مستقل داشته باشد.

\item \textbf{\lr{Test 4}}: جداسازی \texttt{\lr{PID Namespace}} با استفاده از \texttt{\lr{CLONE\_NEWPID}}.

\item \textbf{\lr{Test 5}}: استفاده همزمان از چند flag برای جداسازی مختلف namespace ها.
\end{enumerate}

\begin{LTR}
\begin{lstlisting}
// تست جامع Namespace ها

int
main(int argc, char *argv[])
{
  char hostname[64];
  int pid;

  printf("=== Comprehensive Namespace Test ===\n\n");

  // Test 1: PID Namespace
  printf("Test 1: PID Namespace\n");
  printf("  Current namespace next_pid: %d\n", 
         get_pid_namespace());
  
  pid = fork();
  if(pid == 0) {
    printf("  Child PID: %d\n", getpid());
    exit(0);
  } else {
    wait(0);
  }

  // Test 2: UTS Namespace (Hostname)
  printf("Test 2: UTS Namespace\n");
  if(getHostname(hostname, 64) == 0) {
    printf("  Hostname: %s\n", hostname);
  }
  setHostname("container1", 10);

  // Test 3: Unshare UTS Namespace
  printf("Test 3: Unshare UTS Namespace\n");
  pid = fork();
  if(pid == 0) {
    unshare(CLONE_NEWUTS);
    setHostname("isolated-host", 14);
    if(getHostname(hostname, 64) == 0) {
      printf("  Child hostname: %s\n", hostname);
    }
    exit(0);
  } else {
    wait(0);
  }

  printf("\n=== All Tests Complete ===\n");
  exit(0);
}
\end{lstlisting}
\end{LTR}

نکات مهم این برنامه:
\begin{itemize}
\item استفاده از تابع \texttt{\lr{unshare}} برای جداسازی namespace های والد و فرزند.
\item بررسی \texttt{\lr{CLONE\_NEWPID}} و \texttt{\lr{CLONE\_NEWUTS}} flag ها برای ایجاد namespace های جدید.
\item نمایش اینکه فرزند می‌تواند تغییرات خود را بدون تأثیر بر والد انجام دهد.
\item آزمایش همزمان چند namespace برای درک عملکرد کامل سیستم.
\end{itemize}


\end{document}


\title{گزارش فاز اول پروژه\\
پیاده‌سازی Swap Memory در xv6}
\author{}
\date{}

\begin{document}
\maketitle

\section{مقدمه}
سیستم‌عامل xv6 به‌صورت پیش‌فرض تنها از حافظه‌ی اصلی (RAM) برای نگهداری صفحات استفاده می‌کند و در صورت پر شدن حافظه، تخصیص حافظه با شکست مواجه می‌شود. هدف فاز اول این پروژه، پیاده‌سازی مکانیزم \lr{Swap Memory} است تا سیستم‌عامل بتواند در شرایط کمبود حافظه، صفحات کم‌استفاده را به دیسک منتقل کرده و در صورت نیاز مجدداً آن‌ها را به حافظه بازگرداند.

در این فاز، تمرکز اصلی بر طراحی و پیاده‌سازی سه بخش زیر است:
\begin{itemize}
  \item ایجاد فرآیند کرنلی برای مدیریت Swap
  \item پیاده‌سازی Swap-out (خارج کردن صفحه از حافظه)
  \item پیاده‌سازی Swap-in (بازگرداندن صفحه به حافظه)
\end{itemize}

\section{نمای کلی طراحی}
ایده‌ی اصلی این فاز مشابه سیستم‌های عامل واقعی است:
\begin{itemize}
  \item حافظه‌ی اصلی نقش یک کش را دارد
  \item دیسک نقش فضای پشتیبان \lr{(Backing Store)} را ایفا می‌کند
  \item جدول صفحات مشخص می‌کند که هر صفحه در حال حاضر در RAM است یا روی دیسک
\end{itemize}

به دلیل محدودیت‌های xv6 و برای جلوگیری از پیچیدگی‌های کرنل، عملیات ورودی/خروجی فایل (خواندن و نوشتن فایل swap) در \textbf{فضای کاربر} انجام می‌شود و کرنل تنها هماهنگی و مدیریت را بر عهده دارد.

\section{فرآیند کرنلی Swap (swapd)}
برای مدیریت عملیات Swap-out یک فرآیند کرنلی دائمی (daemon) به نام \texttt{swapd} ایجاد شده است.

\subsection{ایجاد فرآیند کرنلی}
در فایل‌های \texttt{proc.h} و \texttt{proc.c} ساختار \texttt{proc} توسعه داده شده و دو فیلد جدید اضافه شده است:
\begin{lstlisting}
int is_kproc;
void (*kentry)(void);
\end{lstlisting}

این فیلدها مشخص می‌کنند که یک فرآیند، کرنلی است و نقطه‌ی شروع اجرای آن کجاست.

تابع \texttt{create\_kernel\_process} یک فرآیند جدید ایجاد کرده، آن را در صف زمان‌بند قرار می‌دهد و باعث می‌شود که به جای بازگشت به فضای کاربر، وارد تابع کرنلی مشخص‌شده شود.

\subsection{تابع swapd}
تابع \texttt{swapd} به‌صورت یک حلقه‌ی بی‌نهایت اجرا می‌شود:
\begin{itemize}
  \item اگر درخواستی برای Swap وجود نداشته باشد، به خواب می‌رود
  \item در صورت وجود درخواست، یک صفحه را از حافظه خارج می‌کند
  \item پس از آزاد شدن حافظه، فرآیندهای منتظر را بیدار می‌کند
\end{itemize}

\section{شروع Swap-out}
Swap-out زمانی آغاز می‌شود که تخصیص حافظه با شکست مواجه شود.

\subsection{نقطه‌ی شروع در uvmalloc}
در فایل \texttt{vm.c}، در تابع \texttt{uvmalloc}، زمانی که \texttt{kalloc()} حافظه‌ای برنگرداند، فرآیند فراخواننده وارد حالت انتظار می‌شود:
\begin{lstlisting}
while((mem = kalloc()) == 0){
  swap_wait_for_free_page();
}
\end{lstlisting}

این کار باعث می‌شود فرآیند به جای شکست خوردن، منتظر آزاد شدن یک صفحه بماند.

\section{صف درخواست‌های Swap}
در فایل \texttt{swap.c} یک صف حلقوی از درخواست‌ها پیاده‌سازی شده است. همچنین از دو کانال برای هماهنگی استفاده می‌شود:
\begin{itemize}
  \item کانال اول برای خوابیدن \texttt{swapd} در صورت خالی بودن صف
  \item کانال دوم برای خوابیدن فرآیندهایی که منتظر حافظه هستند
\end{itemize}

مکانیزم \texttt{sleep} و \texttt{wakeup} در xv6 تضمین می‌کند که این هماهنگی بدون شرایط رقابتی انجام شود.

\section{انتخاب صفحه‌ی قربانی (LRU)}
برای انتخاب صفحه‌ی مناسب جهت Swap-out، از یک تقریب الگوریتم \lr{LRU} استفاده شده است.

\subsection{استفاده از بیت Accessed}
در معماری RISC-V، سخت‌افزار هنگام دسترسی به یک صفحه، بیت \texttt{PTE\_A} را فعال می‌کند. الگوریتم پیاده‌سازی‌شده به صورت دو مرحله‌ای است:
\begin{itemize}
  \item در مرحله‌ی اول، اگر صفحه‌ای \texttt{PTE\_A} فعال داشته باشد، این بیت پاک می‌شود
  \item در مرحله‌ی دوم، اولین صفحه‌ای که بیت \texttt{PTE\_A} آن صفر باشد، به عنوان قربانی انتخاب می‌شود
\end{itemize}

\section{متادیتای Swap و بیت PTE\_SWP}
برای تشخیص تفاوت بین یک صفحه‌ی نامعتبر و یک صفحه‌ی Swap شده، یک بیت نرم‌افزاری جدید به جدول صفحات اضافه شده است:
\begin{lstlisting}
#define PTE_SWP (1L << ...)
\end{lstlisting}

معانی حالت‌ها:
\begin{itemize}
  \item \texttt{PTE\_V=1}: صفحه در حافظه است
  \item \texttt{PTE\_V=0, PTE\_SWP=1}: صفحه Swap شده و قابل بازیابی است
  \item \texttt{PTE\_V=0, PTE\_SWP=0}: دسترسی نامعتبر (خطای سگمنتیشن)
\end{itemize}

\section{نوشتن صفحات Swap در فضای کاربر}
طبق صورت پروژه، عملیات فایل نباید در کرنل انجام شود. به همین دلیل یک برنامه‌ی کاربری به نام \texttt{swapper} نوشته شده است.

\subsection{ارتباط کرنل و فضای کاربر}
کرنل و برنامه‌ی swapper از طریق یک mailbox مشترک و syscallها با هم ارتباط برقرار می‌کنند:
\begin{itemize}
  \item کرنل درخواست Swap را ثبت می‌کند
  \item swapper درخواست را دریافت کرده و فایل \texttt{.swp} را می‌نویسد یا می‌خواند
  \item نتیجه به کرنل بازگردانده می‌شود
\end{itemize}

نام فایل swap به‌صورت:
\[
\texttt{pid\_[L2,L1,L0].swp}
\]
تعریف شده تا هر صفحه به‌صورت یکتا قابل بازیابی باشد.

\section{Swap-in و Page Fault}
Swap-in زمانی اتفاق می‌افتد که یک فرآیند به صفحه‌ای Swap شده دسترسی پیدا کند.

\subsection{مدیریت Page Fault}
در تابع \texttt{usertrap} در فایل \texttt{trap.c}، هنگام وقوع page fault:
\begin{itemize}
  \item آدرس خطادار استخراج می‌شود
  \item اگر \texttt{PTE\_SWP} فعال باشد، فرآیند Swap-in آغاز می‌شود
\end{itemize}

\subsection{تابع swapin\_page}
مراحل Swap-in:
\begin{enumerate}
  \item تخصیص یک صفحه‌ی فیزیکی جدید
  \item خواندن داده‌ها از فایل swap توسط swapper
  \item بازسازی ورودی جدول صفحات
  \item پاک‌سازی بیت \texttt{PTE\_SWP} و فعال‌سازی \texttt{PTE\_V}
  \item اجرای \texttt{sfence\_vma} برای همگام‌سازی TLB
\end{enumerate}

پس از این مراحل، دستور خطادار مجدداً اجرا شده و برنامه بدون خطا ادامه می‌یابد.

\section{تست فاز اول}
برای بررسی صحت پیاده‌سازی، برنامه‌ای به نام \texttt{swaptest} نوشته شده است:
\begin{itemize}
  \item ایجاد ۲۰ فرآیند فرزند
  \item تخصیص ۱۰ صفحه‌ی ۴KB در هر فرزند
  \item پر کردن صفحات با الگوی مشخص
  \item بررسی صحت داده‌ها پس از Swap-out و Swap-in
\end{itemize}

برای فعال شدن سریع‌تر Swap، مقدار \texttt{PHYSTOP} در فایل \texttt{memlayout.h} کاهش داده شده است.

نمایش پیام \texttt{swaptest: ALL OK} نشان‌دهنده‌ی موفقیت کامل فاز اول است.

\section{نتیجه‌گیری}
در این فاز، یک سیستم Swap کامل شامل Swap-out و Swap-in به xv6 اضافه شد. این پیاده‌سازی:
\begin{itemize}
  \item از نظر معماری به سیستم‌های عامل واقعی نزدیک است
  \item جداسازی صحیح مسئولیت‌ها بین کرنل و فضای کاربر را رعایت می‌کند
  \item تحت فشار حافظه به‌درستی عمل می‌کند
\end{itemize}

با موفقیت تست‌ها، فاز اول پروژه به‌طور کامل تکمیل شده است.

\end{document}
